\section*{第 \ref{ch:g2:pushes-and-pulls} 章——推和拉：力}

\begin{solution}{exo:g2:pushes-and-pulls:1}
推；拉；推；拉；推（挤压就是从两侧同时推）。
\end{solution}

\begin{solution}{exo:g2:pushes-and-pulls:2}
力作用在物体上的位置、力的作用方向，
以及力的大小（用箭头的长短表示）。
\end{solution}

\begin{solution}{exo:g2:pushes-and-pulls:3}
使物体开始运动、停止运动、转向、改变形状。例子各不相同，比如：
弹出一颗弹珠（开始运动）、接住飞盘（停止运动），
让球撞墙后反弹（转向）、压扁一个雪球（改变形状）。
\end{solution}

\begin{solution}{exo:g2:pushes-and-pulls:4}
让它停下：手指逆着弹珠的运动方向推；
让它加速：沿它滚动的方向从后面推；
让它转弯：从侧面横着推，垂直于它的路径。
\end{solution}

\begin{solution}{exo:g2:pushes-and-pulls:5}
能恢复原状的：海绵和橡皮筋。能保持凹陷的：
橡皮泥和松软的雪。
\end{solution}

\begin{solution}{exo:g2:pushes-and-pulls:6}
轻推：一支指向右方的短箭头。重推：
一支指向左方的长箭头。大小不同，方向相反。
\end{solution}

\begin{solution}{exo:g2:pushes-and-pulls:7}
地球的拉力。它不需要接触物体——它隔着\omterm{def:g2:air-around-us:air}{空气}拉着勺子，
并且始终竖直向下。
\end{solution}

\begin{solution}{exo:g2:pushes-and-pulls:8}
地球的拉力（掉落的袜子会下落），以及\omterm{def:g1:magnets:magnet}{磁铁}的拉力（
冰箱\omterm{def:g1:magnets:magnet}{磁铁}隔着纸把图画压在冰箱上）。
\end{solution}

\begin{solution}{exo:g2:pushes-and-pulls:9}
绳子上画两支箭头，长度相等，方向相反。当两边的拉力
\emph{大小相等}时，彩带保持静止——
它们相互抵消。一旦某一队拉得更用力，
它的箭头就占了上风，彩带便朝那一方向移动。
\end{solution}

\begin{solution}{exo:g2:pushes-and-pulls:10}
三个位置受到的是同一个力：地球的拉力，箭头竖直向下——
上升途中（使球减速）、最高点（即将把球拉回）、
下落途中（使球加速）。它
哪怕是一瞬间也从不停止拉拽。
\end{solution}
